\begin{theorem}[Continuation morphism composite visible tails result cases]
\label{thm:continuation-morphism-comp-visible-tails-result-cases}
For histories $a,b,c,l,m$, let
$left:\mathsf{ContinuationMorphism}(a,b)$ and
$right:\mathsf{ContinuationMorphism}(b,c)$. If the two displayed tails are
$\hsame(\mathsf{tail}(left),\Eone(l))$ and
$\hsame(\mathsf{tail}(right),\Eone(m))$, then there exists a history $k$ such
that
$$
\hsame(\mathsf{tail}(left\circ right),\Eone(k))\land
\ContR(\Eone(l),m,k).
$$
\end{theorem}
\begin{proof}
Unpack the two continuation morphism records. The closed composite stores the
append of the two visible tails as its displayed tail, which is definitionally
one-headed. Choose the exposed payload append as $k$; the remaining continuation
ledger is the canonical append continuation from $\Eone(l)$ through $m$.
\end{proof}
\leanchecked{BEDC.Derived.CategoryUp.ContinuationMorphism\_comp\_visible\_tails\_result\_cases}

\begin{theorem}[Continuation morphism visible-tail composite result is not empty]
\label{thm:continuation-morphism-visible-tail-composite-not-empty}
For histories $a,b,c,l,m$, let
$left:\mathsf{ContinuationMorphism}(a,b)$ and
$right:\mathsf{ContinuationMorphism}(b,c)$. If the two displayed tails are
$\hsame(\mathsf{tail}(left),\Eone(l))$ and
$\hsame(\mathsf{tail}(right),\Eone(m))$, then the closed composite tail cannot
be internally the same as $\emp$.
\end{theorem}
\begin{proof}
Unpack the two continuation morphism records and substitute the visible-tail
equalities. The closed composite tail is the append of two one-headed tails,
which is itself one-headed. BEDC constructor separation rules out an internal
history sameness proof from this one-headed tail to $\emp$.
\end{proof}
\leanchecked{BEDC.Derived.CategoryUp.ContinuationMorphism\_comp\_visible\_tails\_result\_not\_empty}

\begin{theorem}[Continuation morphism composite tail nonempty cases]
\label{thm:category-continuation-morphism-comp-tail-nonempty-cases}
For histories $a,b,c$, let
$left:\mathsf{ContinuationMorphism}(a,b)$ and
$right:\mathsf{ContinuationMorphism}(b,c)$. If the closed composite tail is
not internally empty, then at least one displayed factor tail is not internally
empty:
$$
\neg\HSame(\mathsf{tail}(left\circ right),\emp)
\Rightarrow
\neg\HSame(\mathsf{tail}(left),\emp)\vee
\neg\HSame(\mathsf{tail}(right),\emp).
$$
\end{theorem}
\begin{proof}
Unpack the two continuation morphism records and inspect the left displayed
tail. If it is empty, the composite tail reduces to the right displayed tail,
so right-tail emptiness would contradict the composite nonemptiness
hypothesis. If it is zero-headed or one-headed, BEDC constructor separation
rules out any equality from that left tail to $\emp$.
\end{proof}
\leanchecked{BEDC.Derived.CategoryUp.ContinuationMorphism\_comp\_tail\_nonempty\_cases}

\begin{theorem}[Continuation morphism composite tail nonempty endpoint cases]
\label{thm:continuation-morphism-comp-tail-nonempty-endpoint-cases}
For histories $a,b,c$, let
$left:\mathsf{ContinuationMorphism}(a,b)$ and
$right:\mathsf{ContinuationMorphism}(b,c)$. If the closed composite tail is
not internally empty, then at least one adjacent endpoint comparison fails:
$$
\neg\HSame(\mathsf{tail}(left\circ right),\emp)
\Rightarrow
\neg\HSame(a,b)\vee\neg\HSame(b,c).
$$
\end{theorem}
\begin{proof}
First split composite tail nonemptiness into nonemptiness of one displayed
factor tail. For each factor, use the empty-tail endpoint iff for continuation
morphisms: an endpoint comparison would force the corresponding tail to be
internally empty. This contradicts the selected tail-nonempty branch.
\end{proof}
\leanchecked{BEDC.Derived.CategoryUp.ContinuationMorphism\_comp\_tail\_nonempty\_endpoint\_cases}

\begin{theorem}[Continuation morphism composite tail nonempty equivalence]
\label{thm:continuation-morphism-comp-tail-nonempty-iff}
For histories $a,b,c$, let
$left:\mathsf{ContinuationMorphism}(a,b)$ and
$right:\mathsf{ContinuationMorphism}(b,c)$. The closed composite tail is not
internally empty exactly when at least one displayed factor tail is not
internally empty:
$$
\neg\HSame(\mathsf{tail}(left\circ right),\emp)
\Longleftrightarrow
\neg\HSame(\mathsf{tail}(left),\emp)\vee
\neg\HSame(\mathsf{tail}(right),\emp).
$$
\end{theorem}
\begin{proof}
Unpack the two continuation morphism records. The composite tail is the append
of the displayed factor tails, so the append nonempty factor equivalence gives
the displayed statement.
\end{proof}
\leanchecked{BEDC.Derived.CategoryUp.ContinuationMorphism\_comp\_tail\_nonempty\_iff}

\begin{theorem}[Continuation morphism composite target empty iff]
\label{thm:continuation-morphism-comp-target-empty-iff}
For histories $a,b,c$, let
$left:\mathsf{ContinuationMorphism}(a,b)$ and
$right:\mathsf{ContinuationMorphism}(b,c)$. The final target is empty exactly
when both sources and both displayed factor tails are empty:
$$
\HSame(c,\emp)\Longleftrightarrow
\HSame(a,\emp)\land\HSame(b,\emp)\land
\HSame(\mathsf{tail}(left),\emp)\land
\HSame(\mathsf{tail}(right),\emp).
$$
\end{theorem}
\begin{proof}
If $c\hsame\emp$, transport the right continuation ledger to empty target and
apply two-stage empty-target inversion. This yields empty source, empty
middle, and empty tails for the two displayed factors.

Conversely, destruct the right continuation morphism. The empty middle and
empty right-tail witnesses reduce its continuation ledger to the canonical
empty continuation, forcing the target to be $\emp$.
\end{proof}
\leanchecked{BEDC.Derived.CategoryUp.ContinuationMorphism\_comp\_target\_empty\_iff}

\begin{theorem}[Continuation morphism composite nonempty tail rejects empty target]
\label{thm:category-continuation-morphism-comp-tail-nonempty-target-not-empty}
For histories $a,b,c$, let
$left:\mathsf{ContinuationMorphism}(a,b)$ and
$right:\mathsf{ContinuationMorphism}(b,c)$. If the closed composite tail is
not internally empty, then the final target cannot be internally empty:
$$
\neg\HSame(\mathsf{tail}(left\circ right),\emp)
\Rightarrow
\neg\HSame(c,\emp).
$$
\end{theorem}
\begin{proof}
Transport the right continuation relation across an internal equality
$c\hsame\emp$. The empty-target inversion for the two-stage continuation then
forces both displayed tails to be internally empty. The composite-tail
nonempty equivalence turns the composite nonemptiness hypothesis into
nonemptiness of one displayed factor tail, contradicting the inversion data.
\end{proof}
\leanchecked{BEDC.Derived.CategoryUp.ContinuationMorphism\_comp\_tail\_nonempty\_target\_not\_empty}

\begin{theorem}[Continuation morphism composite unary tails reject zero-headed result]
\label{thm:category-continuation-morphism-comp-unary-tail-e0-absurd}
For histories $a,b,c,z$, let
$left:\mathsf{ContinuationMorphism}(a,b)$ and
$right:\mathsf{ContinuationMorphism}(b,c)$. If both displayed tails are unary
histories, then the closed composite tail cannot be internally the same as
$\Ezero(z)$.
\end{theorem}
\begin{proof}
Unpack the two continuation morphism records. The closed composite tail is the
append of the two component tails, and unary closure under append keeps that
composite tail unary. A unary history cannot be internally the same as a
zero-headed history, by the BEDC unary zero-exclusion rule.
\end{proof}
\leanchecked{BEDC.Derived.CategoryUp.ContinuationMorphism\_comp\_unary\_tail\_e0\_absurd}

\begin{theorem}[Continuation morphism composite one-headed tail left cases]
\label{thm:continuation-morphism-comp-e1-tail-left-cases}
For histories $a,b,c,k$, let
$left:\mathsf{ContinuationMorphism}(a,b)$ and
$right:\mathsf{ContinuationMorphism}(b,c)$. If the composite displayed tail is
classified by $\Eone(k)$ and the left displayed tail is unary, then either the
left displayed tail is empty and the right displayed tail is classified by
$\Eone(k)$, or the left displayed tail is itself one-headed and forms a
continuation witness with the right displayed tail to $\Eone(k)$.
\end{theorem}
\begin{proof}
Unpack the two continuation morphism records and inspect the left tail. The
empty case reduces the composite tail by left-unit append. The zero-headed case
contradicts unary classification. In the one-headed case, the displayed
composite equality is exactly the continuation witness from the exposed left
tail and the right tail to $\Eone(k)$.
\end{proof}
\leanchecked{BEDC.Derived.CategoryUp.ContinuationMorphism\_comp\_e1\_tail\_left\_cases}

\begin{theorem}[Continuation morphism composite one-headed tail right cases]
\label{thm:category-continuation-composite-e1-tail-right-cases}
For histories $a,b,c,k$, let
$left:\mathsf{ContinuationMorphism}(a,b)$ and
$right:\mathsf{ContinuationMorphism}(b,c)$. If the composite displayed tail is
classified by $\Eone(k)$ and the right displayed tail is unary, then either the
right displayed tail is empty and the left displayed tail is classified by
$\Eone(k)$, or the right displayed tail is itself one-headed and forms a
continuation witness from the left displayed tail to $\Eone(k)$.
\end{theorem}
\begin{proof}
Unpack the two continuation morphism records and inspect the right tail. The
empty case reduces the composite tail by right-unit append. The zero-headed
case contradicts unary classification. In the one-headed case, the displayed
composite equality is exactly the continuation witness from the left tail and
the exposed right tail to $\Eone(k)$.
\end{proof}
\leanchecked{BEDC.Derived.CategoryUp.ContinuationMorphism\_comp\_e1\_tail\_right\_cases}

\begin{theorem}[Continuation morphism composite zero-headed tail left cases]
\label{thm:continuation-morphism-comp-e0-tail-left-cases}
For histories $a,b,c,k$, let
$left:\mathsf{ContinuationMorphism}(a,b)$ and
$right:\mathsf{ContinuationMorphism}(b,c)$. If the composite displayed tail is
classified by $\Ezero(k)$ and the left displayed tail is unary, then either the
left displayed tail is empty and the right displayed tail is classified by
$\Ezero(k)$, or the left displayed tail is one-headed and forms a continuation
witness with the right displayed tail to $\Ezero(k)$.
\end{theorem}
\begin{proof}
Unpack the two continuation morphism records and inspect the left tail. The
empty case reduces the composite tail by left-unit append. The zero-headed
case contradicts unary classification of the left tail. In the one-headed
case, the displayed composite equality is exactly the continuation witness
from the exposed left tail and the right tail to $\Ezero(k)$.
\end{proof}
\leanchecked{BEDC.Derived.CategoryUp.ContinuationMorphism\_comp\_e0\_tail\_left\_cases}

\begin{theorem}[Continuation morphism composite zero-headed tail right cases]
\label{thm:continuation-morphism-comp-e0-tail-right-cases}
For histories $a,b,c,k$, let
$left:\mathsf{ContinuationMorphism}(a,b)$ and
$right:\mathsf{ContinuationMorphism}(b,c)$. If the composite displayed tail is
classified by $\Ezero(k)$ and the right displayed tail is unary, then the right
displayed tail is empty and the left displayed tail is classified by
$\Ezero(k)$.
\end{theorem}
\begin{proof}
Unpack the two continuation morphism records and inspect the right tail. The
empty case reduces the composite tail by right-unit append. The zero-headed
case contradicts unary classification of the right tail. The one-headed case
would classify a one-headed composite tail as zero-headed, contradicting BEDC
constructor separation.
\end{proof}
\leanchecked{BEDC.Derived.CategoryUp.ContinuationMorphism\_comp\_e0\_tail\_right\_cases}

\begin{theorem}[Continuation morphism nonempty one-headed composite factors]
\label{thm:continuation-morphism-comp-nonempty-e1-tail-factors}
For histories $a,b,c,k$, let
$left:\mathsf{ContinuationMorphism}(a,b)$ and
$right:\mathsf{ContinuationMorphism}(b,c)$. If both displayed tails are unary
and nonempty, and the composite displayed tail is classified by $\Eone(k)$,
then there are histories $l,m$ such that the left and right displayed tails are
classified by $\Eone(l)$ and $\Eone(m)$, and
$$
\ContR(\Eone(l),\Eone(m),\Eone(k)).
$$
\end{theorem}
\begin{proof}
Destructure both displayed tails. Empty tails contradict the corresponding
nonempty hypotheses, and zero-headed tails contradict unary classification.
The remaining one-headed--one-headed case exposes $l$ and $m$ directly, while
the composite displayed equality supplies the continuation ledger.
\end{proof}
\leanchecked{BEDC.Derived.CategoryUp.ContinuationMorphism\_comp\_nonempty\_e1\_tail\_factors}

\begin{theorem}[Continuation morphism nonempty one-headed composite factors are unique]
\label{thm:continuation-morphism-comp-nonempty-e1-tail-factors-unique}
For histories $a,b,c,k$, let
$left:\mathsf{ContinuationMorphism}(a,b)$ and
$right:\mathsf{ContinuationMorphism}(b,c)$. If both displayed tails are unary
and nonempty, and the composite displayed tail is classified by $\Eone(k)$,
then there are histories $l,m$ such that the left and right displayed tails are
classified by $\Eone(l)$ and $\Eone(m)$, the visible boundary satisfies
$$
\ContR(\Eone(l),\Eone(m),\Eone(k)),
$$
and any other visible payloads $l',m'$ with the same two displayed tail
classifications and the same visible boundary satisfy
$\HSame(l,l')\land\HSame(m,m')$.
\end{theorem}
\begin{proof}
Destructure both displayed tails. Empty tails contradict nonemptiness and
zero-headed tails contradict unary classification. In the remaining
one-headed--one-headed case, the witnesses are the two exposed payloads.
Constructor injectivity for $\Eone$ turns any competing left and right tail
classifications into the two payload equalities, while the displayed composite
equality supplies the boundary continuation.
\end{proof}
\leanchecked{BEDC.Derived.CategoryUp.ContinuationMorphism\_comp\_nonempty\_e1\_tail\_factors\_unique}

\begin{theorem}[Continuation morphism right one-headed tail target cases]
\label{thm:continuation-morphism-right-e1-tail-target-cases}
For histories $b,c,m$, let
$right:\mathsf{ContinuationMorphism}(b,c)$. If the displayed tail of $right$
is internally the same as $\Eone(m)$, then the target is one-headed:
$$
\exists r:\mathsf{BHist}.\;
c=\Eone(r)\land \ContR(b,m,r).
$$
\end{theorem}
\begin{proof}
Unpack the continuation morphism record and replace its displayed tail by
$\Eone(m)$. The BEDC continuation step-result inversion then exposes the
unique one-headed target payload $r$ and the descended continuation
$\ContR(b,m,r)$.
\end{proof}
\leanchecked{BEDC.Derived.CategoryUp.ContinuationMorphism\_right\_e1\_tail\_target\_cases}

\begin{theorem}[Continuation morphism right zero-headed tail target cases]
\label{thm:continuation-morphism-right-e0-tail-target-cases}
For histories $b,c,m$, let
$right:\mathsf{ContinuationMorphism}(b,c)$. If the displayed tail of $right$
is internally the same as $\Ezero(m)$, then the target is zero-headed:
$$
\exists r:\mathsf{BHist}.\;
c=\Ezero(r)\land \ContR(b,m,r).
$$
\end{theorem}
\begin{proof}
Unpack the continuation morphism record and replace its displayed tail by
$\Ezero(m)$. The BEDC continuation step-result inversion then exposes the
zero-headed target payload $r$ together with the descended continuation
$\ContR(b,m,r)$.
\end{proof}
\leanchecked{BEDC.Derived.CategoryUp.ContinuationMorphism\_right\_e0\_tail\_target\_cases}

\begin{theorem}[Continuation morphism composite right visible tail result cases]
\label{thm:category-continuation-morphism-comp-right-visible-tail-result-cases}
For histories $a,b,c,m$, let
$left:\mathsf{ContinuationMorphism}(a,b)$ and
$right:\mathsf{ContinuationMorphism}(b,c)$. If the right displayed tail is
$\hsame(\mathsf{tail}(right),\Eone(m))$, then the closed composite tail is
one-headed, and its exposed payload is the continuation of the left tail by
$m$:
$$
\exists k:\mathsf{BHist}.\;
\hsame(\mathsf{tail}(left\circ right),\Eone(k))\land
\ContR(\mathsf{tail}(left),m,k).
$$
\end{theorem}
\begin{proof}
Unpack the two continuation morphism records. The closed composite stores the
append of the left tail with the right tail as its displayed tail. Replacing
the right tail by its $\Eone(m)$ presentation makes the composite tail
definitionally $\Eone(\mathsf{append}(\mathsf{tail}(left),m))$. Choose this
payload append as $k$; the remaining continuation ledger is the reflexive
append continuation.
\end{proof}
\leanchecked{BEDC.Derived.CategoryUp.ContinuationMorphism\_comp\_right\_visible\_tail\_result\_cases}

\begin{theorem}[Continuation morphism composite right zero tail result cases]
\label{thm:category-continuation-morphism-comp-right-e0-tail-result-cases}
For histories $a,b,c,m$, let
$left:\mathsf{ContinuationMorphism}(a,b)$ and
$right:\mathsf{ContinuationMorphism}(b,c)$. If the right displayed tail is
$\hsame(\mathsf{tail}(right),\Ezero(m))$, then the closed composite tail is
zero-headed, and its exposed payload is the continuation of the left tail by
$m$:
$$
\exists k:\mathsf{BHist}.\;
\hsame(\mathsf{tail}(left\circ right),\Ezero(k))\land
\ContR(\mathsf{tail}(left),m,k).
$$
\end{theorem}
\begin{proof}
Unpack the two continuation morphism records. The closed composite stores the
append of the left tail with the right tail as its displayed tail. Replacing
the right tail by its $\Ezero(m)$ presentation makes the composite tail
definitionally $\Ezero(\mathsf{append}(\mathsf{tail}(left),m))$. Choose this
payload append as $k$; the remaining continuation ledger is the reflexive
append continuation.
\end{proof}
\leanchecked{BEDC.Derived.CategoryUp.ContinuationMorphism\_comp\_right\_e0\_tail\_result\_cases}

\begin{theorem}[Continuation morphism composite left-tail result cases]
\label{thm:continuation-morphism-comp-left-tail-result-cases}
For histories $a,b,c,l$, let
$left:\mathsf{ContinuationMorphism}(a,b)$ and
$right:\mathsf{ContinuationMorphism}(b,c)$. If the left displayed tail is
classified by $l$, then the closed composite tail is determined by the right
displayed tail shape:
$$
\HSame(\mathsf{tail}(right),\emp)\land
\HSame(\mathsf{tail}(left\circ right),l),
$$
or there exists a history $m$ such that
$$
\HSame(\mathsf{tail}(right),\Ezero(m))\land
\HSame(\mathsf{tail}(left\circ right),\Ezero(\mathsf{append}(l,m))),
$$
or there exists a history $m$ such that
$$
\HSame(\mathsf{tail}(right),\Eone(m))\land
\HSame(\mathsf{tail}(left\circ right),\Eone(\mathsf{append}(l,m))).
$$
\end{theorem}
\begin{proof}
Unpack the two continuation morphism records and replace the left displayed
tail by its classified representative $l$. Splitting the right displayed tail
into empty, zero-headed, and one-headed cases gives the three composite-tail
shapes directly from the definition of append.
\end{proof}
\leanchecked{BEDC.Derived.CategoryUp.ContinuationMorphism\_comp\_left\_tail\_result\_cases}

\begin{theorem}[Continuation morphism right zero tail composite result is not empty]
\label{thm:continuation-morphism-comp-right-e0-tail-result-not-empty}
For histories $a,b,c,m$, let
$left:\mathsf{ContinuationMorphism}(a,b)$ and
$right:\mathsf{ContinuationMorphism}(b,c)$. If the right displayed tail is
$\hsame(\mathsf{tail}(right),\Ezero(m))$, then the closed composite tail cannot
be internally the same as $\emp$.
\end{theorem}
\begin{proof}
Unpack the two continuation morphism records and replace the right tail by its
zero-headed presentation. The closed composite tail is then a zero-headed
history, namely $\Ezero(\mathsf{append}(\mathsf{tail}(left),m))$. BEDC
constructor separation rules out an internal sameness proof from this
zero-headed history to $\emp$.
\end{proof}
\leanchecked{BEDC.Derived.CategoryUp.ContinuationMorphism\_comp\_right\_e0\_tail\_result\_not\_empty}

\begin{theorem}[Category hom-carrier one-headed target and morphism descent]
\label{thm:category-hom-carrier-e1-target-morphism-exactness}
For histories $a,k,r$, a hom-carrier whose target and morphism histories are
both one-headed descends to the corresponding tail hom-carrier:
$$
\mathsf{CategoryHomCarrier}(a,\Eone(r),\Eone(k))
\Longrightarrow
\mathsf{CategoryHomCarrier}(a,r,k).
$$
\end{theorem}
\begin{proof}
Unpack the hom-carrier fields. The source unary witness is unchanged. Unary
inversion removes the visible constructor from both the target carrier and the
morphism carrier, and the right step-rule inversion for $\ContR$ removes the
same visible constructor from the continuation ledger. Repacking these fields
gives $\mathsf{CategoryHomCarrier}(a,r,k)$.
\end{proof}
\leanchecked{BEDC.Derived.CategoryUp.CategoryHomCarrier\_e1\_target\_morphism\_exactness}

\begin{theorem}[Category hom-carrier one-headed target and morphism exactness iff]
\label{thm:category-hom-carrier-e1-target-morphism-exactness-iff}
For histories $a,k,r$, the simultaneous one-headed target and morphism
presentation is exactly the tail hom-carrier:
$$
\mathsf{CategoryHomCarrier}(a,\Eone(r),\Eone(k))
\Longleftrightarrow
\mathsf{CategoryHomCarrier}(a,r,k).
$$
\end{theorem}
\begin{proof}
The forward implication is the one-headed target and morphism descent theorem.
Conversely, a tail hom-carrier supplies the unary source, unary morphism tail,
and tail continuation $\ContR(a,k,r)$. The one-headed morphism target
exactness theorem repacks these data as the displayed one-headed target and
morphism hom-carrier.
\end{proof}
\leanchecked{BEDC.Derived.CategoryUp.CategoryHomCarrier\_e1\_target\_morphism\_exactness\_iff}

\begin{theorem}[Category hom-carrier one-headed source and morphism target exactness]
\label{thm:category-hom-e1-source-e1-morphism-target-iff}
For histories $a,k,t$,
$$
\begin{aligned}
&\mathsf{CategoryHomCarrier}(\Eone(a),t,\Eone(k)) \\
&\quad\Longleftrightarrow\quad
\exists r,\ t=\Eone(r)\land \mathsf{UnaryHistory}(a)\land
\mathsf{UnaryHistory}(k)\land \mathsf{UnaryHistory}(r)\land
\ContR(\Eone(a),k,r).
\end{aligned}
$$
\end{theorem}
\begin{proof}
Unpack the hom-carrier witness. The continuation field forces the target to
have a visible one constructor; removing that constructor gives
$\ContR(\Eone(a),k,r)$. The source, morphism, and target unary witnesses
invert through the same visible constructors. Conversely, the displayed target
equation, the three unary witnesses, and the one-step continuation rule rebuild
the hom-carrier fields.
\end{proof}
\leanchecked{BEDC.Derived.CategoryUp.CategoryHomCarrier\_e1\_source\_e1\_morphism\_target\_iff}
